% ============================================================================
%  Precautionary Saving in Europe --- A Descriptive Overview
%  Wave 1 of the descriptive scene-setter (ideas 1, 4, 5, 6).
%  Compile with:  pdflatex documentation.tex   (run twice for references)
% ============================================================================
\documentclass[11pt,a4paper]{article}

\usepackage[T1]{fontenc}
\usepackage[utf8]{inputenc}
\usepackage{mathpazo}
\linespread{1.05}
\usepackage{microtype}

\usepackage[a4paper,margin=2.6cm]{geometry}
\usepackage{booktabs}
\usepackage{array}
\usepackage{amsmath,amssymb}
\usepackage{graphicx}
\graphicspath{{figures/}{./}}
\usepackage{caption}
\captionsetup{font=small,labelfont=bf}
\usepackage{enumitem}
\setlist{itemsep=2pt,topsep=3pt}
\usepackage[section]{placeins}
\renewcommand{\topfraction}{0.92}
\renewcommand{\bottomfraction}{0.85}
\renewcommand{\textfraction}{0.08}
\renewcommand{\floatpagefraction}{0.75}
\usepackage{xcolor}
\usepackage[hidelinks,colorlinks=true,linkcolor=black!70!blue,
            urlcolor=black!60!blue,citecolor=black]{hyperref}

\newcommand{\code}[1]{\texttt{#1}}
\newcommand{\dataset}[1]{\textsc{#1}}
\newcommand{\fig}[3]{%
  \begin{figure}[htbp]\centering
    \IfFileExists{figures/#2.png}
      {\includegraphics[width=#1\linewidth]{#2}}
      {\fbox{\parbox{0.80\linewidth}{\centering\itshape figure not found ---
        run \code{bash run.sh} to generate it.}}}
    \caption{#3}\label{fig:#2}
  \end{figure}}

\title{\textbf{Precautionary Saving in Europe: A Descriptive Overview}\\[3pt]
       \large Structural vs Cyclical, the Map, Spending, and What Households Hold}
\author{Project documentation --- descriptive extension}
\date{\today}

\begin{document}
\maketitle
\thispagestyle{empty}

\begin{abstract}
\noindent
Before testing \emph{why} the euro-area household saving rate is high, this note
sets the scene \emph{descriptively}, as the supervisors asked. Four pictures: (1)
is the elevated rate a new \emph{structural} norm or a \emph{cyclical} bump? (4) a
map of how saving rates vary across Europe; (5) how spending rotated from goods to
services; and (6) what households actually hold, and how the risk mix has evolved.
Headlines: the saving-rate \emph{trend} has stepped up about $1.7$ pp
($12.9\%\!\to\!14.6\%$) with the cycle now near zero, so the persistence is mostly
structural; the North--South map (Germany ${\sim}20\%$ vs the South ${\sim}11$--$13\%$,
Greece dissaving) is long-standing, not new; real services consumption rebounded
above its 2019 level while goods flattened, reconciling high saving with strong
services; and the risky (equity \& funds) share of household wealth rose about
$+8$ pp over the decade ($29\%\!\to\!37\%$) even as deposits held near $31\%$. All
figures use free data (Eurostat + GISCO; no FRED) from a live run on \today\ and
should be regenerated before citation. It covers all six descriptive ideas:
structural-vs-cyclical, the map, goods-vs-services, the asset mix, demographics \&
pensions, and the US discrepancy.
\end{abstract}

\tableofcontents
\bigskip\hrule\bigskip

% ============================================================================
\section{Why descriptive, and what these pictures show}\label{sec:intro}
% ============================================================================
The companion \code{../extension\_feedback/} folder asks \emph{why} the saving
rate is high (a precautionary, liquidity-tiered reading). The supervisors asked
to step back first and describe the phenomenon. This note does that with six
self-contained charts, each a standalone script reusing the project helpers and
the house style (shaded COVID / energy-hiking bands, a dashed Feb-2022 line,
direct annotations, an italic footnote). It motivates, rather than tests, the
hypothesis work.

% ============================================================================
\section{Structural or cyclical?}\label{sec:struct}
% ============================================================================
Is the elevated rate a new \emph{norm} or a \emph{cycle} that should revert? A
Hodrick--Prescott filter ($\lambda=1600$) splits the saving rate into a slow trend
(structural) and a cycle (Figure~\ref{fig:structural_vs_cyclical}). The structural
trend has \textbf{stepped up about $1.7$ pp}, from $12.9\%$ (2012--19) to $14.6\%$
today, while the latest cycle is essentially \emph{zero} --- the rate sits at its
trend. The 2020 lockdown spike was almost entirely cyclical and has unwound. So
the persistence is now \textbf{mostly structural}: a higher norm, consistent with
saving more for old age as pension promises weaken, rather than a transient bump.
(The 2020 spike distorts the HP trend around 2020; read the pre-2020 vs latest
step, not the 2020 bend.)

\fig{0.92}{structural_vs_cyclical}{The euro-area saving rate split into a
structural trend and a cyclical component (HP, $\lambda=1600$); the lower panel
shows the cycle against cyclical drivers (ECB rate, inflation, GPR), $z$-scored.}

% ============================================================================
\section{The European map}\label{sec:map}
% ============================================================================
Where is saving high? A choropleth of the gross household saving rate by country
(Figure~\ref{fig:europe_saving_map}): \textbf{Germany ${\sim}20\%$ and France
${\sim}18\%$}, the Nordics high; \textbf{Italy ${\sim}11\%$, Spain ${\sim}13\%$},
and \textbf{Greece dissaving}. The North--South gap is \emph{structural and
long-standing} --- it predates the energy shock --- so a raw "who saves most" map
mostly reflects habits, not the recent crisis.

\fig{0.78}{europe_saving_map}{Gross household saving rate by country (latest year);
Eurostat \dataset{tec00131} with GISCO boundaries. North dark, South light.}

% ============================================================================
\section{Spending: goods vs services}\label{sec:goods}
% ============================================================================
If saving is so high, how has services spending stayed so strong? Real consumption
volumes (2019$=$100, Figure~\ref{fig:goods_vs_services}) resolve the puzzle:
\textbf{services collapsed in 2020 then rebounded to ${\sim}107$}, while goods
boomed in lockdown and then flattened (${\sim}103$). Households kept consuming
services (reopening, and services inflation lifts the \emph{nominal} figure
further) and saved by \textbf{pulling back on goods}. High saving and strong
services are two sides of a \textbf{goods$\rightarrow$services rotation}, not a
contradiction.

\fig{0.86}{goods_vs_services}{Real household consumption of goods vs services
(2019$=$100) against the saving rate; Eurostat \dataset{nama\_10\_fcs}.}

% ============================================================================
\section{What households save in, and how it evolved}\label{sec:compose}
% ============================================================================
Disaggregating the household balance sheet by \emph{risk}
(Figure~\ref{fig:saving_composition_evolution}): the \textbf{risky share (equity \&
investment-fund shares, F5) rose about $+8$ pp over the decade, $29\%\to37\%$},
while \textbf{non-risky deposits (F2) held near $31\%$} and insurance/pensions also
rose. Equity is now the \emph{largest single instrument}, though deposits and
pensions (non-risky) still dominate combined. So the non-risky base is still large,
but \textbf{equity involvement has clearly grown} --- more euro-area households now
carry market exposure. The cross-country snapshot
(Figure~\ref{fig:risky_share_by_country}) shows that risky share is higher in the
core/North and lower in the South.

\fig{0.86}{saving_composition_evolution}{Non-risky (deposits, F2) vs risky (equity
\& funds, F5) share of euro-area household financial assets over time; Eurostat
\dataset{nasa\_10\_f\_bs}.}

\fig{0.74}{risky_share_by_country}{Risky (equity \& investment-fund) share of
household financial assets, by country (latest).}

% ============================================================================
\section{Demographics and pensions}\label{sec:demo}
% ============================================================================
Do aging and pension design move household saving? Two cross-country scatters
(Figure~\ref{fig:demographics_pensions}). \emph{(a)} The old-age dependency ratio
is, if anything, \emph{weakly negatively} related to the saving rate
($r\approx-0.14$): older Southern economies save \emph{less}, not more, so aging
does not mechanically lift saving. \emph{(b)} A funded-pension proxy --- the
insurance \& pension share of household assets (F6), high where households
actually \emph{hold} pension assets and low under pay-as-you-go state promises ---
is \emph{weakly positively} related to saving ($r\approx+0.34$): big funded systems
(Netherlands, Denmark) coexist with solid saving, not low saving. So "good pensions
$\Rightarrow$ less private saving" is not clear in the cross-section --- pensions
shape \emph{where} wealth sits (funded systems hold more in F6/F5) more than
\emph{how much} is saved.

\fig{0.96}{demographics_pensions}{Cross-country: old-age dependency vs the saving
rate (left) and the funded-pension proxy (F6 share) vs the saving rate (right),
latest year; Eurostat \dataset{demo\_pjanind}, \dataset{tec00131},
\dataset{nasa\_10\_f\_bs}.}

% ============================================================================
\section{Why Europeans look like cautious savers}\label{sec:caution}
% ============================================================================
The clearest US--Europe difference is not how \emph{much} households save but
\emph{how} they hold it (Figure~\ref{fig:us_vs_ea_caution}). US households hold
about \textbf{57\%} of their financial wealth in equity \& investment funds versus
\textbf{${\sim}37\%$} in the euro area --- roughly $1.5\times$ --- with the mirror
image in deposits (US ${\sim}11\%$ vs euro area ${\sim}31\%$). The gap is large and
\emph{persistent} across the last two decades. The institutional history behind it:
the US built a funded, equity-heavy private-pension system (401(k)/IRA) and an
equity-investing culture, while much of continental Europe leaned on
\emph{pay-as-you-go state pensions} and a bank-deposit culture. European households
therefore carry less market risk and look more "cautious" --- a structural,
pension-rooted difference, not merely preferences.

\fig{0.86}{us_vs_ea_caution}{Risky (equity \& funds, F5) and non-risky (deposits,
F2) shares of household financial assets, US vs euro area, over time. US: OECD
financial accounts; euro area: Eurostat \dataset{nasa\_10\_f\_bs}.}

% ============================================================================
\section{Synthesis}\label{sec:synth}
% ============================================================================
\begin{enumerate}[label=\arabic*.]
  \item The elevated saving rate is now a \textbf{higher structural norm}, not a
  passing cycle (\S\ref{sec:struct}).
  \item It sits on a \textbf{structural North--South map} that long predates the
  shock (\S\ref{sec:map}).
  \item High saving coexists with \textbf{strong services} via a
  goods$\rightarrow$services rotation (\S\ref{sec:goods}).
  \item Households still hold a large \textbf{non-risky} base, but \textbf{equity
  involvement has grown} markedly over the decade (\S\ref{sec:compose}).
  \item \textbf{Aging and pension design} move household saving only weakly in the
  cross-section; pensions shape \emph{where} wealth sits more than how much is saved
  (\S\ref{sec:demo}).
  \item Versus the US, euro-area households are not low savers so much as
  \textbf{cautious holders} --- far less equity, far more cash --- a structural,
  pension-rooted difference (\S\ref{sec:caution}).
\end{enumerate}
These descriptive facts motivate the hypothesis work in
\code{../extension\_feedback/}: a precautionary, liquidity-tiered reading of why
the rate is high and where the money sits. All series are pulled live and revised
--- regenerate before citing.

\end{document}
